\begin{abstract}
This thesis develops rigorous stochastic partial differential equation (SPDE) descriptions for the dynamics of mean–field particle systems. We analyze two representative SPDEs that characterizes the mesoscopic or macroscopic behaviors (fluctuation and condensation respectively) emerging from microscopic mean-field type stochastic interactions, namely the Dean--Kawasaki (DK) equation and the Anderson Nonlinear Schrödinger equation (NLS).\\

The empirical density of Brownian particles with mean–field interactions in both drift and noise is formally linked to the Dean–Kawasaki equation, a singular conservative SPDE that captures hydrodynamic fluctuations. We introduce a suitable regularization, prove well-posedness in analytic function spaces, and establish its effectiveness as an approximation of the particle system. In particular, we derive weak error estimates and a restricted large-deviation principle, thus validating a stochastic hydrodynamic description;\\

In the mean–field regime, we also consider the Bosonic particles in a random medium driven by a spatial white noise potential. We show that the condensate is governed by the Anderson nonlinear Schrödinger equation (ANLS). Using recent constructions of the continuum Anderson Hamiltonian, we prove local and global well-posedness for ANLS under various initial data and interaction assumptions. We further obtain qualitative and quantitative control of the condensate and convergence to the effective dynamics, providing a rigorous framework for dynamical Bose–Einstein condensation in random environments.
\end{abstract}

\begin{abstractgermain}
   Diese Dissertation entwickelt eine rigorose Beschreibung der Dynamik von Mean-Field-Teilchensystemen mittels stochastischer partieller Differentialgleichungen (SPDEs). Wir untersuchen zwei repräsentative SPDEs, welche das makroskopische Fluktuations- bzw. Kondensationsverhalten beschreiben, das aus mikroskopischen stochastischen Wechselwirkungen hervorgeht, nämlich die Dean–Kawasaki-Gleichung und die Anderson-nichtlineare Schrödinger-Gleichung (NLS).\\

Die empirische Dichte von brownschen Teilchen mit Mean-Field-Wechselwirkungen sowohl im Drift als auch im Rauschen ist formal mit der Dean–Kawasaki-Gleichung verknüpft, einer singulären konservativen SPDE, die hydrodynamische Fluktuationen erfasst. Wir führen eine geeignete Regularisierung ein, beweisen die Wohlgestelltheit in analytischen Funktionsräumen und zeigen die Effektivität dieser Gleichung als Approximation des Teilchensystems. Insbesondere leiten wir schwache Fehlerabschätzungen sowie ein eingeschränktes Großabweichungsprinzip her und validieren damit eine stochastische hydrodynamische Beschreibung.\\

Im Mean-Field-Regime betrachten wir außerdem Bosonen in einem zufälligen Medium, das durch ein räumlich weißes Rauschpotential modelliert wird. Wir zeigen, dass das Kondensat durch die Anderson-nichtlineare Schrödinger-Gleichung (ANLS) beschrieben wird. Unter Verwendung kürzlicher Konstruktionen des Kontinuums--Anderson--Hamilton\-oper\-ators beweisen wir lokale und globale Wohlgestelltheit der ANLS unter verschiedenen Anfangsdaten und Annahmen an die Wechselwirkung. Darüber hinaus erhalten wir qualitative und quantitative Kontrolle des Kondensats sowie die Konvergenz zur effektiven Dynamik und liefern damit einen rigorosen Rahmen für die dynamische Bose–Einstein-Kondensation in zufälligen Umgebungen.
\end{abstractgermain}

\chapter*{Eigenständigkeitserklärung}
\addcontentsline{toc}{chapter}{Eigenständigkeitserklärung}

Hiermit erkläre ich, dass ich die vorliegende Dissertation selbstständig verfasst und keine anderen als die angegebenen Hilfsmittel, Hilfen und Quellen verwendet habe. \\

Zusätzlich erkläre ich, dass Künstliche Intelligenz (KI) in dieser Arbeit ausschließlich als technisches Hilfsmittel zur Erstellung von Abbildungen in der Einleitung verwendet wurde - dies umfasst sowohl Strukturdiagramme als auch Grafiken zur Illustration von Konzepten. Diese rein grafische Unterstützung ist akademisch vollständig unabhängig von den eigentlichen wissenschaftlichen Beiträgen dieser Arbeit. Alle mathematischen Resultate, Herleitungen und Beweise stellen ausnahmslos meine eigene, originäre geistige Leistung dar.

\vspace{2cm}
\noindent
Berlin, den \today \hfill \rule{6cm}{0.4pt}\\
\hspace*{0pt}\hfill Unterschrift